%--------------------------
\begin{frame}{Question 3}

\begin{block}{}
L’administration chargée de l’examen préliminaire international :

\begin{itemize}
  \item[$\circ$] a. a toujours le droit d’inviter le déposant à fournir une copie du document de priorité. 
  \item[$\bullet$] b. a toujours le droit de soulever un défaut d’unité de l’invention. 
  \item[$\bullet$] c. n’a pas le droit de déclarer comme n’ayant pas été présentée une revendication de 
  priorité si elle considère que la validité de cette dernière est en jeu. 
  \item[$\circ$] d. Aucune des réponses a, b ou c n’est correcte.
\end{itemize}

\end{block}

\end{frame}

%--------------------------
\begin{frame}{Question 3}

\textbf{Base juridique :} \pctart{34}(3)(a) + \pctrule{17} + \pctrule{66}.7.  

\begin{itemize}
  \item[$\circ$] \textbf{a.} \; Faux : la copie du document de priorité relève en principe de l’office récepteur
  ou du Bureau international
  (\pctrule{17}.1).

  \item[$\bullet$] \textbf{b.} \; Correct : l’IPEA peut toujours examiner l’unité de l’invention et inviter
  au paiement de taxes additionnelles
  (\pctart{34}.3.a)).

  \item[$\bullet$] \textbf{c.} \; Correct : l’IPEA ne peut pas déclarer une priorité comme non présentée
  pour des raisons de validité matérielle ; elle peut seulement en tenir compte dans son opinion
  (\pctrule{66}.7).

  \item[$\circ$] \textbf{d.} \; Faux puisque \textbf{b} et \textbf{c} sont correctes.
\end{itemize}

\end{frame}